\begin{proof}[Proof of \cref{obj:lem:rectangular-sampling-bound}]
\begin{enumerate}
\item Define the dimension-dependent constants
\[
M_d=4^{2d},\qquad
D_d=q4^{2d},\qquad
C_{\mathrm{ker},d}=M_d^4(D_d^2+1),
\]
\[
C_{\mathrm{coord},d}=6M_d^2+12C_{\mathrm{ker},d},
\qquad
C=(q+q^2+1)C_{\mathrm{coord},d}.
\]
These constants are finite and positive. Fix an admitted law \(P\) and admissible \(n,m,h,J\), and set
\[
V_{n,m,h,J}
=\frac{1}{nh^d}+\frac{k}{nNh^{2d}}\ge0.
\]
We will bound each response coordinate and each raw-matrix coordinate by
\(C_{\mathrm{coord},d}V_{n,m,h,J}\), then sum the bounds after symmetrization.
% lean: CausalSmith.Stat.TwosamplePointcateAnnotationFrontier.rectangular_sampling_bound

\item We first establish the explicit basis bounds used in the coordinate estimates. Write
\[
\mathcal U_d
=\left\{u\in\mathbb N^d:\sum_{i=1}^d u_i\le2\right\},
\qquad |\mathcal U_d|=q,
\]
and define the reference polynomial factors, for \(\xi_{\mathrm{ref}}\in\mathbb R\), by
\[
\begin{aligned}
p_{\mathrm{ref},0}(\xi_{\mathrm{ref}})&=1,\\
p_{\mathrm{ref},1}(\xi_{\mathrm{ref}})&=\sqrt{12}\,\xi_{\mathrm{ref}},\\
p_{\mathrm{ref},2}(\xi_{\mathrm{ref}})
&=\sqrt5\left(6\xi_{\mathrm{ref}}^2-\frac12\right).
\end{aligned}
\]
For \(|\xi_{\mathrm{ref}}|\le1/2\),
\[
0\le\xi_{\mathrm{ref}}^2\le\frac14,
\qquad
\left|6\xi_{\mathrm{ref}}^2-\frac12\right|\le1,
\]
and consequently
\[
|p_{\mathrm{ref},0}(\xi_{\mathrm{ref}})|=1\le4,\qquad
|p_{\mathrm{ref},1}(\xi_{\mathrm{ref}})|
\le\frac{\sqrt{12}}2\le4,\qquad
|p_{\mathrm{ref},2}(\xi_{\mathrm{ref}})|\le\sqrt5\le3\le4.
\]
The product formula in \cref{obj:synth_1} therefore gives
\[
|r_u(x)|\le4^d
\qquad(x\in C_h,\ u\in\mathcal U_d).
\]

For completeness, the same factors give an explicit row-integral bound for the fine kernel. Set
\[
\delta=\frac hJ,\qquad
\mathcal Z_J=\{0,\ldots,J-1\}^d,\qquad
x_{\mathrm{cell},\mathbf z,i}
=\frac12-\frac h2+\left(z_i+\frac12\right)\delta
\quad(\mathbf z\in\mathcal Z_J).
\]
Let \(C_{\mathrm{cell},\mathbf z}\) be the product of the intervals
\[
\left[\frac12-\frac h2+z_i\delta,\,
      \frac12-\frac h2+(z_i+1)\delta\right),
\]
with the right endpoint included when \(z_i=J-1\). Thus these cells partition \(C_h\).
The fine basis coordinates of \cref{obj:def:projection} are
\[
b_{\mathbf z,u}(x)
=\sqrt{J^d}\,\mathbf1_{C_{\mathrm{cell},\mathbf z}}(x)
 \prod_{i=1}^d
 p_{\mathrm{ref},u_i}
 \left(\frac{x_i-x_{\mathrm{cell},\mathbf z,i}}{\delta}\right),
\qquad
x\in\mathbb R^d,\quad
(\mathbf z,u)\in\mathcal Z_J\times\mathcal U_d.
\]
In particular, their values are zero outside their supporting cells, and
\[
|b_{\mathbf z,u}(x)|\le\sqrt{J^d}\,4^d.
\]
With \(\boldsymbol0=(0,\ldots,0)\), the constant coordinate satisfies
\[
b_{\mathbf z,\boldsymbol0}(x)
=\sqrt{J^d}\,\mathbf1_{C_{\mathrm{cell},\mathbf z}}(x),
\qquad
\int b_{\mathbf z,\boldsymbol0}(x)^2\,d\nu_h(x)=1,
\]
where the last identity is its orthonormal normalization. Pointwise,
\[
|b_{\mathbf z,u}(x)|
\le \frac{4^d}{\sqrt{J^d}}\,
      b_{\mathbf z,\boldsymbol0}(x)^2,
\]
so integration yields
\[
\int |b_{\mathbf z,u}(x)|\,d\nu_h(x)
\le\frac{4^d}{\sqrt{J^d}}.
\]
If \(x\in C_{\mathrm{cell},\mathbf z}\), only that cell contributes to the kernel, and hence
\[
K_J(x,x')
=\sum_{u\in\mathcal U_d}
 b_{\mathbf z,u}(x)b_{\mathbf z,u}(x').
\]
It follows that
\[
\begin{aligned}
\int |K_J(x,x')|\,d\nu_h(x')
&\le\sum_{u\in\mathcal U_d}
 |b_{\mathbf z,u}(x)|
 \int |b_{\mathbf z,u}(x')|\,d\nu_h(x')\\
&\le q\bigl(\sqrt{J^d}\,4^d\bigr)
          \frac{4^d}{\sqrt{J^d}}
=D_d.
\end{aligned}
\]
If \(x\) belongs to no fine cell, every fine basis coordinate at \(x\) is zero, and the same bound holds. Thus
\[
\int |K_J(x,x')|\,d\nu_h(x')\le D_d
\qquad(x\in\mathbb R^d).
\]
The basis-product formula also gives
\[
K_J(x,x')=K_J(x',x).
\]
% lean: CausalSmith.Stat.TwosamplePointcateAnnotationFrontier.fineKernel_abs_row_bound

\item We now express the two coordinate types through a common marked kernel. Define the role counts
\[
\ell=\lfloor n/2\rfloor,\qquad t=N-\ell.
\]
Both are positive because \(n\ge2\) and \(N\ge n\). Let \(W\) and \(Z\) be independent generic treatment-role and outcome-role records, respectively:
\[
W=(X^T,A^T)\in[0,1]^d\times\{0,1\},
\qquad
Z=(X^L,A^L,Y^L)\in[0,1]^d\times\{0,1\}^2,
\]
\[
\mathcal L(W)=P_{XA},
\qquad
\mathcal L(Z)=\mathcal L_P(X,A,Y).
\]
Their covariate marginals are uniform by \cref{obj:def:primitive-class}. The product sampling law and the disjoint roles in \cref{obj:def:roles} give independent records
\[
W_i=(X_i^T,A_i^T)\quad(i\in I_T),
\qquad
Z_j=(X_j,A_j,Y_j)\quad(j\in I_L),
\]
independent within and between the two roles.

Use the coordinate tags
\[
\diamond\in
\bigl(\{\mathrm R\}\times\mathcal U_d\bigr)
\sqcup
\bigl(\{\mathrm Q\}\times\mathcal U_d\times\mathcal U_d\bigr).
\]
For each tag, define the three mark functions by
\[
\begin{array}{c|ccc}
\diamond
& f_{\diamond,0}(Z)
& f_{\diamond,T}(W)
& f_{\diamond,L}(Z)\\ \hline
(\mathrm R,u)
& r_u(X^L)A^LY^L
& r_u(X^T)A^T
& Y^L\\
(\mathrm Q,u,v)
& r_u(X^L)r_v(X^L)A^L
& r_u(X^T)A^T
& r_v(X^L)A^L
\end{array}
\]
on the full record domains displayed above. Binary marks have absolute value at most one. The coarse-coordinate bound from the preceding step therefore gives
\[
\begin{aligned}
|f_{\diamond,0}(Z)|&\le M_d &&(X^L\in C_h),\\
|f_{\diamond,T}(W)|&\le M_d &&(X^T\in C_h),\\
|f_{\diamond,L}(Z)|&\le M_d &&(X^L\in C_h).
\end{aligned}
\]
Define
\[
H_\diamond(W,Z)
=w_h(X^T)w_h(X^L)K_J(X^T,X^L)
 f_{\diamond,T}(W)f_{\diamond,L}(Z),
\]
\[
\widehat S_\diamond
=\frac1\ell\sum_{j\in I_L}
 w_h(X_j)f_{\diamond,0}(Z_j),
\qquad
\widehat H_\diamond
=\frac1{t\ell}\sum_{i\in I_T}\sum_{j\in I_L}
 H_\diamond(W_i,Z_j),
\qquad
F_\diamond=\widehat S_\diamond-\widehat H_\diamond.
\]
These expressions vanish termwise whenever a localization weight is zero. The rectangular formulas in \cref{obj:def:moments} give exactly
\[
F_{(\mathrm R,u)}=(\widehat R_{h,J})_u,
\qquad
F_{(\mathrm Q,u,v)}
=(\widehat Q^{\mathrm{raw}}_{h,J})_{uv}.
\]
In the response identity, expanding the empirical projection gives the rectangular formula on \(C_h\); outside \(C_h\), its defining zero value and the zero localization weight give the same identity.
% lean: CausalSmith.Stat.TwosamplePointcateAnnotationFrontier.rectangular_r_coordinate_bound

\item The localized mark and kernel energies admit explicit bounds. Since
\[
w_h(x)=h^{-d}\mathbf1_{C_h}(x),
\qquad
\mathbb P(X^L\in C_h)=\mathbb P(X^T\in C_h)=h^d,
\]
uniform design gives
\[
\mathbb E\,w_h(X^L)^2
=\mathbb E\,w_h(X^T)^2=h^{-d}.
\]
Thus
\[
\mathbb E\!\left[
 \bigl(w_h(X^L)f_{\diamond,0}(Z)\bigr)^2
 \right]\le \frac{M_d^2}{h^d}.
\]
Similarly, bounding both marks and integrating over the two independent covariates gives
\[
\begin{aligned}
\mathbb E\,H_\diamond(W,Z)^2
&\le M_d^4h^{-4d}
 \iint_{C_h\times C_h}K_J(x,x')^2\,dx\,dx'\\
&=\frac{M_d^4}{h^{2d}}
 \iint K_J(x,x')^2\,d\nu_h(x)\,d\nu_h(x')\\
&=\frac{M_d^4k}{h^{2d}},
\end{aligned}
\]
using the kernel-square identity in
\cref{obj:lem:local-polynomial-projection-facts}.

Define the integrated slices on the full record domains by
\[
H_{\diamond,T}(W)=\mathbb E_Z H_\diamond(W,Z),
\qquad
H_{\diamond,L}(Z)=\mathbb E_W H_\diamond(W,Z),
\]
where each expectation integrates an independent record with the law specified above. The row-integral estimate gives
\[
\begin{aligned}
\left|
 \mathbb E_Z\!\left[
 w_h(X^L)K_J(x,X^L)f_{\diamond,L}(Z)
 \right]\right|
&\le M_d\,\mathbb E_Z\!
 \left[w_h(X^L)|K_J(x,X^L)|\right]\\
&=M_d\int |K_J(x,x')|\,d\nu_h(x')\\
&\le M_dD_d.
\end{aligned}
\]
Consequently,
\[
|H_{\diamond,T}(W)|
\le M_d^2D_d\,w_h(X^T).
\]
Kernel symmetry gives the corresponding bound with the roles reversed:
\[
|H_{\diamond,L}(Z)|
\le M_d^2D_d\,w_h(X^L).
\]
Squaring and integrating yields
\[
\mathbb E\,H_{\diamond,T}(W)^2
\le\frac{M_d^4D_d^2}{h^d},
\qquad
\mathbb E\,H_{\diamond,L}(Z)^2
\le\frac{M_d^4D_d^2}{h^d}.
\]
All the displayed features and kernels are measurable and square-integrable. In particular, their means, centered squares, and finite-sum integrals are well defined.
% lean: CausalSmith.Stat.TwosamplePointcateAnnotationFrontier.rectangular_marked_kernel_row_energy

\item We derive the rectangular variance formula used with these energies. Define
\[
\begin{aligned}
\mu_\diamond&=\mathbb E\,H_\diamond(W,Z),\\
H_{\diamond,0}(W,Z)&=H_\diamond(W,Z)-\mu_\diamond,\\
H_{\diamond,1}(W)&=H_{\diamond,T}(W)-\mu_\diamond,\\
H_{\diamond,2}(Z)&=H_{\diamond,L}(Z)-\mu_\diamond,\\
H_{\diamond,12}(W,Z)
&=H_{\diamond,0}(W,Z)-H_{\diamond,1}(W)-H_{\diamond,2}(Z).
\end{aligned}
\]
Fubini's theorem gives
\[
\mathbb E H_{\diamond,0}
=\mathbb E H_{\diamond,1}
=\mathbb E H_{\diamond,2}=0,
\]
and
\[
\begin{aligned}
\mathbb E\!\left[H_{\diamond,0}H_{\diamond,1}(W)\right]
&=\mathbb E H_{\diamond,1}(W)^2,\\
\mathbb E\!\left[H_{\diamond,0}H_{\diamond,2}(Z)\right]
&=\mathbb E H_{\diamond,2}(Z)^2,\\
\mathbb E\!\left[H_{\diamond,1}(W)H_{\diamond,2}(Z)\right]
&=0.
\end{aligned}
\]
Indeed, the first two identities follow by integrating the other record first, and the third follows from independence and centering. Expanding the residual square therefore gives
\[
\mathbb E H_{\diamond,12}^2
=\mathbb E H_{\diamond,0}^2
 -\mathbb E H_{\diamond,1}^2
 -\mathbb E H_{\diamond,2}^2.
\]
Also,
\[
\mathbb E H_{\diamond,0}^2
=\mathbb E H_\diamond^2-\mu_\diamond^2
\le\mathbb E H_\diamond^2.
\]
Hence
\[
\mathbb E H_{\diamond,12}^2\le\mathbb E H_\diamond^2,
\]
and centering each slice gives
\[
\mathbb E H_{\diamond,1}^2
\le\mathbb E H_{\diamond,T}^2,
\qquad
\mathbb E H_{\diamond,2}^2
\le\mathbb E H_{\diamond,L}^2.
\]

To compute the rectangular second moment, recall the independent-sample identity. For \(n_{\mathrm{av}}\ge1\) independent identically distributed square-integrable real variables \(Z_{\mathrm{av},j}\), define
\[
\overline Z_{\mathrm{av}}
=\frac1{n_{\mathrm{av}}}
 \sum_{j=1}^{n_{\mathrm{av}}}Z_{\mathrm{av},j}.
\]
Expanding the square and factoring the mixed expectations gives
\[
\mathbb E\,\overline Z_{\mathrm{av}}^2
=\frac1{n_{\mathrm{av}}}\mathbb E Z_{\mathrm{av},1}^2
 +\left(1-\frac1{n_{\mathrm{av}}}\right)
   (\mathbb E Z_{\mathrm{av},1})^2.
\]
Apply this identity first to the outcome-record average, with the treatment records fixed, and then to the treatment-record averages. Finite-sum integration commutes with taking the slice means. Since \(H_{\diamond,0}\) has mean zero, the resulting identity is
\[
\begin{aligned}
\mathbb E(\widehat H_\diamond-\mu_\diamond)^2
&=\frac{\mathbb E H_{\diamond,0}^2}{t\ell}\\
&\quad+\frac1t\left(1-\frac1\ell\right)
       \mathbb E H_{\diamond,1}^2\\
&\quad+\left(1-\frac1t\right)\frac1\ell
       \mathbb E H_{\diamond,2}^2.
\end{aligned}
\]
Substituting the residual-energy identity above reduces this to
\[
\mathbb E(\widehat H_\diamond-\mu_\diamond)^2
=\frac{\mathbb E H_{\diamond,1}^2}{t}
 +\frac{\mathbb E H_{\diamond,2}^2}{\ell}
 +\frac{\mathbb E H_{\diamond,12}^2}{t\ell}.
\]
The three preceding energy comparisons imply
\[
\mathbb E(\widehat H_\diamond-\mu_\diamond)^2
\le
\frac{\mathbb E H_{\diamond,T}^2}{t}
+\frac{\mathbb E H_{\diamond,L}^2}{\ell}
+\frac{\mathbb E H_\diamond^2}{t\ell}.
\]
% lean: CausalSmith.Stat.TwosamplePointcateAnnotationFrontier.rectangular_ustat_variance

\item We convert this inequality to the sampling envelope and then include the first-order term. The floor split satisfies
\[
\ell\ge1,\qquad n\le2\ell+1\le3\ell,
\qquad 2\ell\le n\le N.
\]
Therefore
\[
\ell\ge\frac n3,\qquad
t=N-\ell\ge\frac N2,
\]
and
\[
\frac1\ell\le\frac3n,\qquad
\frac1t\le\frac2N\le\frac2n,\qquad
\frac1{t\ell}\le\frac6{nN}.
\]
These inequalities apply for every \(m\ge0\), including \(m=0\).

By the definition of \(C_{\mathrm{ker},d}\),
\[
M_d^4D_d^2\le C_{\mathrm{ker},d},
\qquad
M_d^4\le C_{\mathrm{ker},d}.
\]
The slice and kernel bounds thus give
\[
\begin{aligned}
\mathbb E(\widehat H_\diamond-\mu_\diamond)^2
&\le \frac{5C_{\mathrm{ker},d}}{nh^d}
     +\frac{6C_{\mathrm{ker},d}k}{nNh^{2d}}\\
&\le6C_{\mathrm{ker},d}V_{n,m,h,J}.
\end{aligned}
\]
The independent-sample mean identity also gives
\[
\mathbb E\widehat S_\diamond
=\mathbb E\!\left[w_h(X^L)f_{\diamond,0}(Z)\right],
\qquad
\mathbb E\widehat H_\diamond=\mu_\diamond,
\]
and
\[
\begin{aligned}
\mathbb E(\widehat S_\diamond-\mathbb E\widehat S_\diamond)^2
&=\frac1\ell
 \left\{
 \mathbb E\!\left[
   \bigl(w_h(X^L)f_{\diamond,0}(Z)\bigr)^2
 \right]
 -\left(\mathbb E[w_h(X^L)f_{\diamond,0}(Z)]\right)^2
 \right\}\\
&\le\frac{3M_d^2}{nh^d}.
\end{aligned}
\]
Linearity of expectation gives the centered-difference identity
\[
F_\diamond-\mathbb EF_\diamond
=(\widehat S_\diamond-\mathbb E\widehat S_\diamond)
 -(\widehat H_\diamond-\mu_\diamond).
\]
The pointwise inequality
\[
(\xi_{\mathrm{diff}}-\zeta_{\mathrm{diff}})^2
\le2\xi_{\mathrm{diff}}^2+2\zeta_{\mathrm{diff}}^2
\qquad(\xi_{\mathrm{diff}},\zeta_{\mathrm{diff}}\in\mathbb R)
\]
follows from \((\xi_{\mathrm{diff}}+\zeta_{\mathrm{diff}})^2\ge0\). Applying it to the centered difference yields
\[
\begin{aligned}
\mathbb E(F_\diamond-\mathbb EF_\diamond)^2
&\le\frac{6M_d^2}{nh^d}
  +12C_{\mathrm{ker},d}V_{n,m,h,J}\\
&\le C_{\mathrm{coord},d}V_{n,m,h,J},
\end{aligned}
\]
because \(1/(nh^d)\le V_{n,m,h,J}\). This argument accommodates the shared outcome records in the two terms.
% lean: CausalSmith.Stat.TwosamplePointcateAnnotationFrontier.rectangular_marked_moment_bound

\item Applying the preceding estimate to the response tags gives
\[
\mathbb E\left[
 \bigl((\widehat R_{h,J})_u-\bar R_u\bigr)^2
 \right]
\le C_{\mathrm{coord},d}V_{n,m,h,J}
\qquad(u\in\mathcal U_d).
\]
Applying it to the raw-matrix tags gives
\[
\mathbb E\left[
 \bigl((\widehat Q^{\mathrm{raw}}_{h,J})_{uv}
       -\mathbb E(\widehat Q^{\mathrm{raw}}_{h,J})_{uv}\bigr)^2
 \right]
\le C_{\mathrm{coord},d}V_{n,m,h,J}
\qquad(u,v\in\mathcal U_d).
\]
The square-integrability established above applies to every coordinate in these two displays.
% lean: CausalSmith.Stat.TwosamplePointcateAnnotationFrontier.rectangular_q_coordinate_bound

\item Define the centered raw matrix by
\[
B_{\mathrm{raw},uv}
=(\widehat Q^{\mathrm{raw}}_{h,J})_{uv}
 -\mathbb E(\widehat Q^{\mathrm{raw}}_{h,J})_{uv},
\qquad u,v\in\mathcal U_d.
\]
The symmetrization formula in \cref{obj:def:moments} and linearity of expectation give
\[
(Q_{P,h,J})_{uv}
=\frac{
 \mathbb E(\widehat Q^{\mathrm{raw}}_{h,J})_{uv}
 +\mathbb E(\widehat Q^{\mathrm{raw}}_{h,J})_{vu}}2,
\]
and therefore
\[
(\widehat Q_{h,J}-Q_{P,h,J})_{uv}
=\frac{B_{\mathrm{raw},uv}+B_{\mathrm{raw},vu}}2.
\]
Every symmetrized coordinate is square-integrable as an average of two square-integrable raw coordinates.

For each sample,
\[
\left(\frac{B_{\mathrm{raw},uv}+B_{\mathrm{raw},vu}}2\right)^2
\le\frac{B_{\mathrm{raw},uv}^2+B_{\mathrm{raw},vu}^2}{2},
\]
since the difference between the right and left sides is
\((B_{\mathrm{raw},uv}-B_{\mathrm{raw},vu})^2/4\). Summing and exchanging the two indices gives the pointwise contraction
\[
\begin{aligned}
\|\widehat Q_{h,J}-Q_{P,h,J}\|_F^2
&\le\frac12\sum_{u,v\in\mathcal U_d}
 \left(B_{\mathrm{raw},uv}^2+B_{\mathrm{raw},vu}^2\right)\\
&=\sum_{u,v\in\mathcal U_d}B_{\mathrm{raw},uv}^2.
\end{aligned}
\]
% lean: CausalSmith.Stat.TwosamplePointcateAnnotationFrontier.rectangular_symmetrization_energy_le

\item Integrability permits exchanging expectation and each finite sum. Summing the response-coordinate bounds gives
\[
\begin{aligned}
\mathbb E\|\widehat R_{h,J}-\bar R\|_2^2
&=\sum_{u\in\mathcal U_d}
 \mathbb E\left[
   \bigl((\widehat R_{h,J})_u-\bar R_u\bigr)^2
 \right]\\
&\le qC_{\mathrm{coord},d}V_{n,m,h,J}.
\end{aligned}
\]
Integrating the pointwise matrix contraction and summing the raw-coordinate bounds gives
\[
\begin{aligned}
\mathbb E\|\widehat Q_{h,J}-Q_{P,h,J}\|_F^2
&\le\sum_{u,v\in\mathcal U_d}
 \mathbb E B_{\mathrm{raw},uv}^2\\
&\le q^2C_{\mathrm{coord},d}V_{n,m,h,J}.
\end{aligned}
\]
Adding these inequalities, and using
\(C_{\mathrm{coord},d}V_{n,m,h,J}\ge0\), yields
\[
\begin{aligned}
&\mathbb E\|\widehat R_{h,J}-\bar R\|_2^2
+\mathbb E\|\widehat Q_{h,J}-Q_{P,h,J}\|_F^2\\
&\qquad\le(q+q^2)C_{\mathrm{coord},d}V_{n,m,h,J}\\
&\qquad\le(q+q^2+1)C_{\mathrm{coord},d}V_{n,m,h,J}\\
&\qquad=C\left\{\frac1{nh^d}+\frac{k}{nNh^{2d}}\right\}.
\end{aligned}
\]
The chosen constant depends on \(d\), and hence has the permitted dependence on the fixed public parameters. All estimates hold uniformly over the admitted law and the admissible sample sizes and tuning inputs.
% lean: CausalSmith.Stat.TwosamplePointcateAnnotationFrontier.rectangular_sampling_bound
\end{enumerate}
\end{proof}
